\documentclass[11pt]{amsart}
\usepackage{amssymb,amsmath}
\usepackage[hidelinks]{hyperref}
\newtheorem{theorem}{Theorem}
\newtheorem{proposition}[theorem]{Proposition}
\newtheorem{lemma}[theorem]{Lemma}
\newtheorem{corollary}[theorem]{Corollary}
\theoremstyle{remark}
\newtheorem{remark}[theorem]{Remark}
\newcommand{\Z}{\mathbb{Z}}
\newcommand{\Q}{\mathbb{Q}}
\newcommand{\F}{\mathbb{F}}
\DeclareMathOperator{\rk}{rank}
\setlength{\emergencystretch}{2em}

\title[The $7$-rank of a putative $\mathrm{srg}(99,14,1,2)$]{On the $7$-rank of a putative strongly regular graph with parameters $(99,14,1,2)$}
\author{Dilan Djematene}
\date{\today}

\begin{document}
\begin{abstract}
Let $A$ be the adjacency matrix of a hypothetical strongly regular graph with parameters $(99,14,1,2)$, the
object of Conway's $99$-graph problem. The only interesting prime for the $p$-ranks of $A$ is $7$, and the known
bound is $\rk_7(A-3I)\le 45$. We prove that $\rk_7(A-3I)$ is odd and that
\[
23\le\rk_7(A-3I)\le 43 .
\]
The upper bound comes from the genus of the integral eigenvalue lattice $\Z^{99}\cap\ker(A+4I)$, through the oddity
formula. The lower bound comes from an exact local formula: the $7$-rank is $15$ plus the $7$-rank of an explicit
$84\times84$ matrix built from the second neighbourhood of any vertex. A small exact computation raises the lower bound
to $25$.
\end{abstract}
\maketitle

\section{Introduction}
A strongly regular graph $\Gamma$ with parameters $(99,14,1,2)$ has spectrum $14^1$, $3^{54}$, $(-4)^{44}$. Its
adjacency matrix satisfies
\begin{equation}\label{eq:quad}
A^2+A-12I=2J,\qquad\text{that is}\qquad (A-3I)(A+4I)=2J .
\end{equation}
Whether such a graph exists is one of Conway's \$1000 problems~\cite{Con}. The $p$-ranks of $A-\theta I$ are classical
invariants of strongly regular graphs~\cite{BvE}. For primes $p$ not dividing $r-s=7$ these ranks are determined by the parameters
(see~\cite{BvE}), so $p=7$ is the only interesting prime here. Brouwer and van Eijl~\cite{BvE} prove
$\rk_p(A-rI)\le\min(f,g)+1$, which gives $\rk_7(A-3I)\le 45$. We write
\[
k_7:=\rk_7(A-3I)-15 .
\]

\begin{theorem}\label{thm:main}
If $\Gamma$ exists, then $k_7$ is even and $8\le k_7\le 28$, that is
\[
\rk_7(A-3I)\in\{23,25,\dots,43\}.
\]
Moreover $k_7=\rk_7(M_x)$ for every vertex $x$, where $M_x$ is the $84\times84$ matrix of Proposition~\ref{prop:local}.
\end{theorem}

The upper bound and the parity are proved in Section~\ref{sec:genus}. The local formula and the lower bound are proved in
Sections~\ref{sec:local} and~\ref{sec:lower}. Remark~\ref{rem:25} explains how a small exhaustive computation excludes
$k_7=8$ as well.

\section{The genus of the eigenvalue lattice}\label{sec:genus}
Let $\Lambda=\Z^{99}\cap\ker(A+4I)$. It is a primitive sublattice of $\Z^{99}$ of rank $44$. The orthogonal projection
onto $\ker(A+4I)$ is
\[
E=\frac{27I-9A+J}{63}.
\]
We use two standard facts.
\begin{enumerate}
\item[(F1)] For a primitive sublattice $L\subseteq\Z^n$, the orthogonal projection of $\Z^n$ onto $L\otimes\Q$ is the
dual lattice $L^\#$.
\item[(F2)] A primitive sublattice $L$ of a unimodular lattice and its orthogonal complement have isomorphic discriminant
groups~\cite[Prop.~1.6.1]{Nik}.
\end{enumerate}

\begin{lemma}\label{lem:disc}
$\Lambda$ is even, and $\Lambda^\#/\Lambda\cong\Z/9\oplus(\Z/7)^{\alpha}$ with $\alpha=\rk_7(A-3I)-1=14+k_7$. In
particular $\det\Lambda=9\cdot7^\alpha$.
\end{lemma}
\begin{proof}
\emph{Evenness.} Every $u\in\Lambda$ is orthogonal to $\mathbf 1$, so $u\cdot u\equiv\sum u_i=0\pmod2$.

\emph{Exponent.} Since $63E$ is integral, (F1) gives $63\Lambda^\#=63E(\Z^{99})\subseteq\Lambda$. So the discriminant
group is killed by $63$, and $\Lambda$ is unimodular at every prime other than $3$ and $7$.

\emph{The $3$-part.} By (F2) it equals the $3$-part of the discriminant group of
$\Lambda'=\Z^{99}\cap(\Q\mathbf 1\oplus\ker(A-3I))$. Write $u\in\Lambda'$ as $u=c\mathbf 1+w$ with $(A-3I)w=0$. Then
$(A-3I)u=11c\mathbf 1$ is integral, so $11c\in\Z$, and $\Lambda'$ contains $\Z\mathbf 1\oplus\Lambda_3$ with index
dividing $11$, where $\Lambda_3=\Z^{99}\cap\ker(A-3I)$. The projection onto $\ker(A-3I)$ is $(44I+11A-2J)/77$. As
above, $77\Lambda_3^\#\subseteq\Lambda_3$, so $\Lambda_3$ is unimodular at $3$. Hence the $3$-part of the discriminant
group of $\Lambda'$ is that of $\Z\mathbf 1$, namely $\Z/9$.

\emph{The $7$-part.} Let $W=\Z^{99}\cap\mathbf 1^\perp$. On $W$ we have $J=0$, hence $A-3I=-7E$ by the formula for $E$.
Since $[\Z^{99}:W\oplus\Z\mathbf 1]=99$ is prime to $7$, we get $\Z_7^{99}=W_7\oplus\Z_7\mathbf 1$ (subscripts denote
$\otimes\Z_7$), and $(A-3I)(W_7)=7E(\Z_7^{99})=7\Lambda^\#_7$. Reducing modulo $7$, the image of $A-3I$ is spanned by
$7\Lambda^\#_7$, which is $\Lambda^\#_7/\Lambda_7$ in $\F_7^{99}$, together with $(A-3I)\mathbf 1=11\cdot\mathbf 1$.
Moreover $\mathbf 1$ is not in the reduction of $7\Lambda^\#_7$, because $\mathbf 1\cdot\mathbf 1=99\not\equiv0$ while
$7\Lambda^\#_7\perp\mathbf 1$. Hence $\rk_7(A-3I)=\dim_{\F_7}(\Lambda^\#/\Lambda)_7+1$. The $7$-part has exponent $7$,
so it is $(\Z/7)^\alpha$ with $\alpha=\rk_7(A-3I)-1$.
\end{proof}

\begin{theorem}\label{thm:genus}
$\alpha$ is even and $\alpha\le42$; equivalently $k_7$ is even and $k_7\le28$. More precisely, write
$\Lambda\otimes\Z_7=U\perp 7V$ with $U,V$ unimodular over $\Z_7$, of ranks $30-k_7$ and $14+k_7$. Then $U\ne0$, and
$U\otimes\F_7$ is not a hyperbolic quadratic space.
\end{theorem}
\begin{proof}
We use the oddity formula of Conway and Sloane~\cite[Ch.~15, \S7]{CS}. For a lattice of signature $\sigma$,
\[
\sigma+\sum_{p\ \mathrm{odd}}p\text{-excess}\equiv\text{oddity}\pmod8 .
\]
For an odd prime $p$, a Jordan constituent of scale $q=p^j$, dimension $n$ and sign $\varepsilon$ contributes
$n(q-1)+4[\,q\text{ is not a square and }\varepsilon=-1\,]$ to the $p$-excess.

Here $\sigma=44$. Since $\Lambda$ is even of odd determinant, it is unimodular and even at $2$, and its oddity is $0$. By
Lemma~\ref{lem:disc} the $3$-adic constituents are a unimodular one and $9^{\pm1}$; the latter contributes
$1\cdot(9-1)=8\equiv0$ whatever its sign, as $9$ is a square. The $7$-adic constituents are $U$ and $7V$, with $\dim V=\alpha$;
the constituent $7V$ contributes $6\alpha+4\delta$, where $\delta=1$ if $\big(\tfrac{\det V}{7}\big)=-1$ and $\delta=0$
otherwise. The oddity formula therefore reads
\[
44+6\alpha+4\delta\equiv0\pmod8 .
\]
Hence $\alpha$ is even; write $\alpha=14+2m$, so that $k_7=2m$. Then $44+84+12m+4\delta\equiv4(m+\delta)\pmod 8$, and so
$\delta\equiv m\pmod2$.

Since $\det\Lambda=9\cdot7^\alpha$, we have $\big(\tfrac{\det U}{7}\big)\big(\tfrac{\det V}{7}\big)=\big(\tfrac{9}{7}\big)=1$.

Suppose first that $U=0$, that is $k_7=30$ and $m=15$. Then $\big(\tfrac{\det V}{7}\big)=1$, so $\delta=0$, contradicting
$\delta\equiv m=15$. Thus $U\ne0$ and $k_7\le28$.

In general, $\big(\tfrac{\det U}{7}\big)=\big(\tfrac{\det V}{7}\big)=(-1)^{m}$. A quadratic space of even dimension $2t=30-k_7$
over $\F_7$ is hyperbolic if and only if its determinant is $(-1)^t$ modulo squares. Here $t=15-m$ and
$\big(\tfrac{-1}{7}\big)=-1$, so hyperbolicity would require $(-1)^m=(-1)^{15-m}$, which never holds.
\end{proof}

\begin{remark}
The same computation applied to graphs that exist is consistent. For the Petersen graph ($r-s=3$), the lattice
$\Z^{10}\cap\ker(A+2I)$ is entirely $3$-divisible in the sense above ($U=0$). This is allowed there, because the unit part
of its determinant, $5$, is a non-square modulo $3$. For $(99,14,1,2)$ the unit part $9$ is a square modulo $7$, and this is
what forbids $U=0$.
\end{remark}

\section{An exact local formula}\label{sec:local}
Fix a vertex $x$. Its neighbourhood consists of $7$ disjoint edges, the \emph{mate pairs}. We write $\bar a$ for the
mate of $a\in N(x)$. Since $\mu=2$, every vertex at distance $2$ from $x$ is adjacent to exactly two vertices of $N(x)$.
These two vertices are not mates, by $\lambda=1$, and every non-mate pair occurs exactly once. So the $84$ vertices at
distance $2$ are labelled by the $84$ non-mate pairs $\{a,b\}\subset N(x)$.

Let $H$ be the adjacency matrix of the second neighbourhood, indexed by these pairs. For two pairs $y,z$ let
$s(y,z)=|y\cap z|$, and let $p(y,z)$ be the number of $(i,j)\in y\times z$ with $j=\bar i$.

\begin{proposition}\label{prop:local}
Let $M_x=H+4I+3S+P'+5J$ over $\F_7$, where $S=(s(y,z))$ and $P'=(p(y,z))$. Then
\[
\rk_7(A-3I)=15+\rk_7(M_x).
\]
In particular $\rk_7(M_x)=k_7$ does not depend on $x$.
\end{proposition}
\begin{proof}
Order the vertices as $\{x\}\cup N(x)$ followed by the second neighbourhood. Then
$A-3I=\begin{pmatrix}F-3I&B\\ B^{T}&H-3I\end{pmatrix}$, where $F$ is the adjacency matrix of the closed neighbourhood
(a windmill), and the column of $B$ at $y=\{a,b\}$ is $e_a+e_b$. A direct computation gives
$\det(F-3I)=2^{23}$, so $F-3I$ is invertible over $\F_7$. Moreover the block of $(F-3I)^{-1}$ on $N(x)$ is
\[
G=-\tfrac38I-\tfrac18P+\tfrac1{16}J ,
\]
where $P$ is the mate involution. To check this, note that the rows of $F-3I$ on $N(x)$ are $(1,\,P-3I)$, and multiply
out using $P^2=I$, $PJ=J$.

So $\rk_7(A-3I)=15+\rk_7(H-3I-B^{T}GB)$. Now $(B^TGB)_{yz}=\sum_{i\in y,\,j\in z}G_{ij}$. Modulo $7$ we have
$-\tfrac38\equiv4$, $-\tfrac18\equiv6$ and $\tfrac1{16}\equiv4$. Since $|y|=|z|=2$, this gives
$(B^TGB)_{yz}\equiv4s+6p+4\cdot4=4s+6p+2$. Hence
\[
H-3I-B^TGB\equiv H-3I+3S+P'+5J\equiv M_x \pmod 7 .\qedhere
\]
\end{proof}

\section{Lower bounds from stars}\label{sec:lower}
For $a\in N(x)$, the \emph{star} of $a$ is the set of the $12$ pairs $\{a,b\}$ with $b\notin\{a,\bar a\}$. Two pairs
$\{a,b\}$ and $\{a,c\}$ are adjacent if and only if they are matched by the perfect matching $\tau_a$ on the $12$ elements
$b$ defined by the triangles through $a$ (the edge $a\{a,b\}$ lies in a unique triangle). The six mate pairs $\{b,\bar b\}$
($b\ne a,\bar a$) are called \emph{classes}.

The \emph{class graph} of $\tau_a$ is the $2$-regular multigraph on the six classes with one edge
$\{\mathrm{class}(b),\mathrm{class}(\tau_a b)\}$ for each pair of $\tau_a$. A loop arises when $\tau_a(b)=\bar b$. Let
$e(\tau_a)$ be the number of cycles of even length in the class graph; a double edge is a cycle of length $2$, and loops
are not counted. Clearly $e(\tau_a)\le3$.

\begin{lemma}\label{lem:star}
The principal submatrix of $M_x$ on the star of $a$ has $7$-rank $11-2e(\tau_a)$.
\end{lemma}
\begin{proof}
On the star, $S=I+J$, and $H+P'=:N$ is the adjacency matrix of the union of the two perfect matchings $\tau_a$ and
$b\mapsto\bar b$. So the block is $N+7I+8J\equiv N+J$.

The union of two perfect matchings is a disjoint union of alternating cycles. A loop of the class graph gives a double
edge $\left(\begin{smallmatrix}0&2\\2&0\end{smallmatrix}\right)$, which is invertible modulo $7$. A cycle of length $m\ge2$
of the class graph gives one alternating cycle of length $2m$. Over $\F_7$, a vector in the kernel of the adjacency matrix
of the cycle $C_{2m}$ satisfies $v_{i+2}=-v_i$. Such vectors exist (forming a space of dimension $2$) exactly when
$4\mid 2m$, that is when $m$ is even, and then they have sum $0$. Hence $\dim\ker_7N=2e(\tau_a)$.

Finally $N\mathbf 1=2\cdot\mathbf 1$, so $(N+J)\mathbf 1=14\cdot\mathbf 1\equiv0$. If $(N+J)v=0$, then
$w=v+\tfrac12(\sum v)\mathbf 1$ lies in $\ker N$. Hence $\ker_7(N+J)=\ker_7N\oplus\F_7\mathbf 1$ has dimension
$2e(\tau_a)+1$, and the rank is $12-2e(\tau_a)-1$.
\end{proof}

\begin{corollary}\label{cor:6}
$k_7\ge 11-2e(\tau_a)$ for every vertex $x$ and every $a\in N(x)$. In particular $k_7\ge6$.
\end{corollary}
\begin{proof}
A principal submatrix has rank at most $\rk_7(M_x)=k_7$. Then use $e\le3$ and the parity of $k_7$
(Theorem~\ref{thm:genus}).
\end{proof}

\begin{proposition}\label{prop:8}
$k_7\ne6$. Hence $k_7\ge8$ and $\rk_7(A-3I)\ge23$.
\end{proposition}
\begin{proof}
Suppose $k_7=6$. By Corollary~\ref{cor:6} every star block has rank at most $6$, hence rank exactly $5$, and
$e(\tau_a)=3$: the class graph of every $\tau_a$ consists of three double edges.

Since $M_x$ is symmetric of rank $6$ over $\F_7$, we can write $M_x=V^{T}QV$ with $V$ of size $6\times84$ and $Q$
invertible. Let $V_a$ be the columns of $V$ indexed by the star of $a$, and $U_a$ their span. Since
$\rk(V_a^TQV_a)=5$, we have $\dim U_a\in\{5,6\}$. If $\dim U_a=6$, then $V_a^TQV_a$ would have rank $6$. Hence
$\dim U_a=5$ and $Q|_{U_a}$ is nondegenerate. Consequently every vector $c$ in the kernel of the star block satisfies
$V_ac=0$, that is,
\[
\sum_{y}c_y\,(M_x)_{y,\cdot}=0\quad\text{on all $84$ columns.}
\]

Take a double edge of the class graph: classes $\{\ell_1,\ell_4\}$ and $\{\ell_2,\ell_3\}$ with
$\tau_a(\ell_1)=\ell_2$, $\tau_a(\ell_3)=\ell_4$, $\bar\ell_2=\ell_3$ and $\bar\ell_4=\ell_1$. In the proof of
Lemma~\ref{lem:star}, the columns $\ell_1$ and $\ell_3$ of $N$ are both $e_{\ell_2}+e_{\ell_4}$. So
$c=e_{\ell_1}-e_{\ell_3}$ is in the kernel of the star block, and the rows of $M_x$ at $y=\{a,\ell_1\}$ and
$y'=\{a,\ell_3\}$ coincide.

Choose $q\in N(x)\setminus\{a,\bar a,\ell_1,\bar\ell_1,\ell_3\}$ (nine choices) and let $z=\{\ell_1,q\}$. By the definition
of $M_x$,
\[
(M_x)_{yz}=H_{yz}+3+5,\qquad (M_x)_{y'z}=H_{y'z}+[\,q=\ell_2\,]+5 .
\]
Here we used $s(y,z)=1$, $p(y,z)=0$ (as $q\ne\bar a$ and $\bar\ell_1\notin z$), $s(y',z)=0$, and
$p(y',z)=[q=\bar\ell_3]=[q=\ell_2]$. Equality modulo $7$ gives $H_{yz}-H_{y'z}\equiv[\,q=\ell_2\,]-3\in\{4,5\}\pmod 7$,
which is impossible since $H_{yz},H_{y'z}\in\{0,1\}$.
\end{proof}

\begin{remark}[$k_7\ne8$]\label{rem:25}
The argument of Proposition~\ref{prop:8} extends. If $k_7=K$, the kernel vectors of a star block that are relations
between complete rows of $M_x$ form a space $R_a$ of dimension at least $12-K$. Each column $z$ outside the star imposes
a linear condition on $R_a$, and this condition depends only on the two (unknown) neighbours of $z$ in the star, so it
ranges over a small finite set. For $K=8$ an exhaustive computation over $\F_7$ shows:
\begin{itemize}
\item for two of the three admissible star types, no subspace of the required dimension meets all the constraint sets;
\item for the third type, exactly $32$ subspaces remain. Adding the degree and labelling constraints of the second
neighbourhood, each of the $32$ resulting Boolean instances is refuted by unit propagation alone.
\end{itemize}
Hence $k_7\ge10$, that is $\rk_7(A-3I)\ge25$. We regard this as computer-assisted, although the computation is tiny. For
$K=10$ the one-star method no longer excludes all types.
\end{remark}

\begin{remark}
A proof that $\rk_7(A-3I)\le 21$, that is $k_7\le6$, would therefore prove that no $\mathrm{srg}(99,14,1,2)$ exists.
\end{remark}

\section*{Acknowledgements}
This work was carried out with the assistance of Claude (Anthropic), an AI system, which was used for derivations,
computations, independent checks of the arguments and the preparation of the manuscript. The author has verified the
results and takes full responsibility for the content.

\begin{thebibliography}{9}
\bibitem{BvE} A.~E.~Brouwer and C.~A.~van Eijl, On the $p$-rank of the adjacency matrices of strongly regular graphs,
\emph{J. Algebraic Combin.} \textbf{1} (1992), 329--346.
\bibitem{Con} J.~H.~Conway, Five \$1{,}000 problems (update 2017), \emph{The On-Line Encyclopedia of Integer Sequences},
\url{https://oeis.org/A248380/a248380.pdf}.
\bibitem{CS} J.~H.~Conway and N.~J.~A.~Sloane, \emph{Sphere Packings, Lattices and Groups}, 3rd ed., Grundlehren der
mathematischen Wissenschaften 290, Springer, New York, 1999.
\bibitem{Nik} V.~V.~Nikulin, Integral symmetric bilinear forms and some of their applications,
\emph{Math. USSR Izv.} \textbf{14} (1980), 103--167.
\end{thebibliography}
\end{document}
